\documentclass[12pt]{article}
\usepackage[utf8]{inputenc}
\usepackage[T1]{fontenc}
\usepackage{hyperref}
\title{Topoiyo Language Resources for K-12 Curriculum Development}
\author{Compiled by Claude Sonnet 5 (claude-sonnet-5) for Andr\'as Kornai}
\date{}
\begin{document}
\maketitle

\section{Introduction}
Topoiyo (Topoyo; ISO 639-3 toy) is an Austronesian language of Topoyo District, Central Mamuju Regency, West Sulawesi, Indonesia. Wikipedia reports about 450 native speakers in 2011. Speakers usually also speak Mamuju and Indonesian. The search returned very little about the language itself.

\section{Primary Educational Resources}
No textbooks, primers or school programmes were found.

\section{Linguistic and Reference Materials}
\begin{itemize}
\item Wikipedia article on Topoiyo. It gives location, speaker count and classification (Kaili--Pamona). \url{https://en.wikipedia.org/wiki/Topoiyo_language}
\item Mamuju language page from the Endangered Languages Alliance, for the neighbouring language and regional context. \url{https://elalliance.org/languages/mamuju}
\end{itemize}
No dictionary or grammar of Topoiyo was found.

\section{Religious and Community Resources}
No scripture was found. Joshua Project describes the Topoyo people as nearly all Muslim. \url{https://joshuaproject.net/people_groups/15562}

\section{Digital and Audio Resources}
No audio, video or software resources were found.

\section{Curriculum Development Context}
Little is documented. The first step would be documentation, with local speakers and linguists in West Sulawesi. Indonesian and Mamuju are the likely languages of schooling and wider communication.
\end{document}